% Generated draft: requires mathematical and source review.
\documentclass[11pt]{article}
\usepackage[margin=25mm]{geometry}
\usepackage{fontspec}
\setmainfont{Latin Modern Roman}
\usepackage{amsmath,amssymb,mathtools,unicode-math}
\setmathfont{Latin Modern Math}
\usepackage{enumitem,needspace,xurl,hyperref}
\setlength{\parindent}{0pt}
\setlength{\parskip}{5pt}
\setlength{\emergencystretch}{3em}
\newcommand{\N}{\mathbb{N}}
\newcommand{\Z}{\mathbb{Z}}
\newcommand{\Q}{\mathbb{Q}}
\newcommand{\R}{\mathbb{R}}
\newcommand{\C}{\mathbb{C}}
\newcommand{\F}{\mathbb{F}}
\newcommand{\A}{\mathbb{A}}
\newcommand{\PP}{\mathbb{P}}
\newcommand{\E}{\mathbb{E}}
\newcommand{\eps}{\varepsilon}
\newcommand{\dd}{\,\mathrm{d}}
\begin{document}
\section*{UnsolvedMath ID: 30001140}
\textbf{Source problem code:} catalogue entry
\subsection{Definitions and mathematical statement}
The supplied record contains two separate questions, retained together under one permanent identifier.

\textbf{(1) Planar development of spherical curves.} Let \(S^2=\{x\in\R^3:\|x\|=1\}\) have its usual spherical metric. A piecewise-geodesic curve is a curve consisting of finitely many geodesic arcs. Simplicity means that the curve has no self-intersections, allowing the identification of the initial and terminal points when the curve is closed. The source does not specify whether the general curve must be closed.

Given a simple piecewise-geodesic curve \(\gamma\) on \(S^2\), consider its development into a planar polygonal chain, preserving the length of each geodesic piece and the corresponding angles at successive vertices. What conditions on \(\gamma\) ensure that this development has no self-intersections? In particular, what can be asserted when \(\gamma\) is a star-shaped spherical polygon?

The phrase ``preserving angles'' requires an angle and orientation convention that the record does not supply: interior angles, unsigned angles between incident arcs, and signed turning angles must be distinguished. The source also does not specify the meaning of spherical star-shapedness, a hemisphere restriction, or whether the development of a closed spherical curve is required to close in the plane. Accordingly, this question does not yet specify a unique development construction or a precise star-shaped subclass.

\textbf{(2) Higher-dimensional discrete integrability.} For an integer \(n\geq4\), consider a proposed \(n\)-dimensional discrete system. The stated consistency requirement is that copies of the system can be imposed simultaneously on every \(n\)-dimensional face of an \((n+1)\)-dimensional cube, compatibly on their intersections. Here these faces are the subsets of \([0,1]^{n+1}\) obtained by fixing one coordinate to \(0\) or \(1\).

Do discrete integrable systems in dimensions four and higher exist with this face-consistency property?

The record supplies neither the unknowns and their locations nor the equations, admissible initial data, nontriviality conditions, or definition of integrability. It also does not state whether consistency requires existence, uniqueness, or independence of the order of extending initial data. Thus compatibility above records only the supplied qualitative requirement; it is not a complete definition of the intended consistency property. The record does not resolve whether existence is requested in every dimension \(n\geq4\) or in some such dimension.
\subsection{Short English statement}
The record asks when a length- and angle-preserving planar development of a simple spherical curve avoids self-intersection, especially for star-shaped polygons. It separately asks for discrete integrable systems in dimensions at least four that are consistent across the faces of a cube one dimension higher. Both questions require additional conventions or hypotheses before they become fully specified mathematical assertions.
\subsection{Sources}
\begin{itemize}
\item Supplied local UnsolvedMath registry record, permanent ID \(30001140\), exact source code \texttt{catalogue entry}, titled \emph{Polygon Development and Higher-Dimensional Discrete Integrability}, in the Oberwolfach Reports Open Problems collection.
\item Supplied report citation: \emph{Discrete Differential Geometry} (2009), Oberwolfach Reports, pp.\ 75--144, DOI \texttt{10.4171/owr/2009/02}. \url{https://doi.org/10.4171/owr/2009/02}. The background describes the report as documenting mathematical workshops and their open-problem discussions. No individual proposer is supplied.
\item Supplied literature source: Rote, \emph{Open Problems in Discrete Differential Geometry}. \url{https://page.mi.fu-berlin.de/rote/Kram/OWR-DDG09-problems.pdf}.
\item Supplied literature source labeled ``original OWR report'': \url{https://doi.org/10.4171/owr/2009/02}.
\item Literature assessment dated exactly \texttt{2026-08-21}: the supplied curation labels the record open and describes both questions as open-form research questions. This is supplied dated curation derived from the OWR material, not an independently verified literature assessment; time-sensitive status and cited claims require independent verification.
\item The supplied statement status is \texttt{corrected\_verified}, with the supplied explanation that both geometric questions are explicit, report-preface prose was removed, and dimensions were normalized. These are registry annotations, not independent verification of the formulation or mathematical status.
\end{itemize}
\end{document}
